%% ════════════════════════════════════════════════════════════════════════════
\section{General-Exponent Carreau-Yasuda: Why Amplitude Perturbation Theory
         Fails, and What Survives}
\label{sec:cy_firstprinciples}
%% ════════════════════════════════════════════════════════════════════════════

Section~\ref{sec:carreau} derives a slow amplitude equation for the classical
Carreau law, $\mu_{\mathrm{eff}}=\mu_0[1+(\lambda_c\dot\gamma)^2]^{(n-1)/2}$,
valid for $\mathrm{Oh}\ll 1$ and small oscillation amplitude. That derivation
is correct for the model it treats. This section asks a harder question:
does an analogous small-amplitude expansion exist for the \emph{general}
Carreau-Yasuda law,
\begin{equation}
  \eta(\dot\gamma) = \eta_\infty + (\eta_0-\eta_\infty)\bigl[1+(\lambda_c\dot\gamma)^a\bigr]^{(n-1)/a},
  \label{eq:cy_general}
\end{equation}
with the shape exponent $a$ left general rather than fixed at $2$ -- the form
this repo's own fitted validation fluid actually uses
($\mathrm{Oh}_0\approx 57.37$, $\lambda_c\approx 3.05\times 10^4$,
$a\approx 0.743$, $\varepsilon_{ST}\equiv(1-n)/a\approx 0.9996$)?

The short answer, reached only after an extensive derivation and two
independent adversarial checks of every step, is: \textbf{no, not for this
fluid's fitted parameters}. This is not a failure of nerve or a shortcut --
it is a precisely quantified mathematical fact about where this particular
fluid's physics actually lives on the Carreau-Yasuda curve. What survives
the investigation is three genuinely new, general, CAS-verified results that
have nothing to do with any specific fluid's parameters, plus a clear,
quantitative diagnosis of exactly what a first-principles treatment of
\emph{this} fluid requires instead.

%% ────────────────────────────────────────────────────────────────────────────
\subsection{Why the classical $O(\epsilon^3)$ route is not available here}
\label{sec:cy_why_not_epsilon3}
%% ────────────────────────────────────────────────────────────────────────────

Section~\ref{sec:carreau}'s expansion rests on one fact: with $a=2$ fixed,
$(\lambda_c\dot\gamma)^2$ is an even integer power of $\dot\gamma$, hence an
analytic (Taylor-expandable) function of the oscillation amplitude
$\epsilon$, since $\dot\gamma=O(\epsilon)$ linearly in Reid's single-mode
ansatz. The viscosity correction is then $O(\epsilon^2)$, the stress
correction $O(\epsilon^3)$, and a standard Landau-equation-style
method-of-averaging calculation goes through cleanly.

\begin{notebox}[The exponent matters, not just its size]
For a \emph{general} exponent $a$, $(\lambda_c\dot\gamma)^a\sim
(\lambda_c\epsilon)^a\,\hat g(\theta)^a$ for some $O(1)$ angular shape
$\hat g$. If $a$ is not an even integer, $|\epsilon|^a$ is not differentiable
at $\epsilon=0$ (it is only H\"older continuous there). For this fluid's
fitted $a\approx 0.743 <1$, the viscosity correction scales as
$|\epsilon|^{0.743}$, which is \emph{larger} than $\epsilon$ itself for small
$\epsilon$ (since $|\epsilon|^{a-1}\to\infty$ as $\epsilon\to0$ when $a<1$).
There is no regime, however small the amplitude, in which this correction is
a small perturbation of a fixed integer-power hierarchy: it sits at a
genuinely fractional order between Reid's $O(\epsilon)$ linear terms and the
usual $O(\epsilon^2)$ geometric mode-coupling correction.
\end{notebox}

This alone does not doom a perturbative treatment -- a fractional-order
expansion parameter is unusual but not unheard of. What actually closes off
the small-amplitude route for this fluid is quantitative, not structural,
and is established in the next section.

%% ────────────────────────────────────────────────────────────────────────────
\subsection{Two natural anchors, and why neither supports a small correction}
\label{sec:cy_anchors}
%% ────────────────────────────────────────────────────────────────────────────

A shear-thinning fluid's viscosity interpolates between two Newtonian
limits: $\eta_0$ at rest ($\dot\gamma\to0$) and $\eta_\infty$ at infinite
shear ($\dot\gamma\to\infty$). Reid's characteristic equation
(Section~\ref{sec:char_eq}) gives the \emph{exact} damping rate
$\lambda_l(\mathrm{Oh})$ and squared frequency $\omega_l^2(\mathrm{Oh})$ at
\emph{any} Ohnesorge number, so either limit is a legitimate place to try to
anchor a correction.

\begin{derivbox}[Anchoring at rest, $\mathrm{Oh}_0$: fails outright]
At $\mathrm{Oh}_0\approx 57.37$, every mode this fluid's solver resolves is
heavily overdamped:
\[
  \begin{array}{cccc}
    l & \lambda_l & \omega_l^2 & \\
    2 & 203.8 & 7.48 & \text{overdamped} \\
    3 & 459.4 & 25.48 & \text{overdamped} \\
    5 & 1119.3 & 102.9 & \text{overdamped} \\
    10 & 3681.4 & 665.4 & \text{overdamped}
  \end{array}
\]
There is no oscillation to average over, so a multiple-scales/envelope
treatment has no fast carrier wave to modulate. Separately, solving for the
shear rate at which the viscosity correction reaches even the $1\%$ level
gives $\dot\gamma\lesssim 7\times10^{-8}$ -- a shear rate no resolved mode
velocity in this solver ever produces. Both the qualitative premise (an
oscillator to perturb) and the quantitative premise (a small correction)
fail simultaneously.
\end{derivbox}

\begin{derivbox}[Anchoring at infinite shear, $\mathrm{Oh}_\infty$: better
                 posed, but still not small]
Define $\mathrm{Oh}_\infty\equiv\mathrm{Oh}_0\,(\eta_\infty/\eta_0) =
\mathrm{Oh}_0(1-\varepsilon_{ST})\approx 0.0254$. Unlike $\mathrm{Oh}_0$,
this \emph{is} comfortably underdamped: at $l=2$,
$\lambda_2=0.116$, $\omega_2^2=7.95$. Reid's oscillator convention is
$\ddot A+2\lambda\dot A+\omega^2A=\text{forcing}$
(\texttt{julia/src/reid.jl}), so the quality factor is
$Q\equiv\sqrt{\omega_2^2}/(2\lambda_2)\approx 12$ -- a genuine separation
between the decay time and the oscillation period, exactly the condition a
weakly-damped multiple-scales treatment needs. A live solver trace (three
impact energies, fine time resolution) confirms the drop's actual effective
Ohnesorge number, once excited, is permanently pinned away from
$\mathrm{Oh}_0$ after a single clean transient at first contact, with no
recurrence on later contact/lift-off cycles.

But ``pinned away from $\mathrm{Oh}_0$'' is not the same claim as ``close to
$\mathrm{Oh}_\infty$.'' Evaluating the exact Carreau-Yasuda law at the shear
rates this fluid's live trace actually visits ($\dot\gamma\sim
0.05$--$1.5$, in the solver's nondimensional units) gives
$\mathrm{Oh}_{\mathrm{eff}}/\mathrm{Oh}_\infty$ ratios of roughly $2$ to
$7$ -- not a small departure from the anchor. And the relevant question is
not this ratio directly, but whether the \emph{first-order Taylor expansion}
of Reid's own exact $\lambda_l(\mathrm{Oh})$, $\omega_l^2(\mathrm{Oh})$
around $\mathrm{Oh}_\infty$ reproduces their true values over that range.
Checked directly against the exact (continuation-solved) relations, at four
points spanning the operative band, $\mathrm{Oh}=0.03,0.05,0.08,0.1$
(ratios $1.2$--$3.9\times\mathrm{Oh}_\infty$): $\omega_l^2$ is nearly flat in
$\mathrm{Oh}$ near $\mathrm{Oh}_\infty$ -- linear-approximation error stays
under $0.3\%$ for $l=2$ across the whole band, and is still small for
higher $l$, growing to about $1.5\%$ by $l=10$ at the band's upper end. But
$\lambda_l$ is far more sensitive and its error \emph{grows across the
band}, not just with $l$: at the low end ($\mathrm{Oh}=0.03$) every $l$
tested is under $0.3\%$, but by the upper end ($\mathrm{Oh}=0.1$) the error
has grown to $8\%$ ($l=2$), $12\%$ ($l=3$), $17\%$ ($l=5$), and $26\%$
($l=10$) -- the same growing-with-$l$ curvature that broke the
$\mathrm{Oh}_0$ anchor reappears here, mode-dependently, and here it also
grows across the operative band itself, not just across modes.
Moving the anchor to some other, empirically ``typical'' $\mathrm{Oh}^\ast$
does not rescue this: the curvature in $\lambda_l(\mathrm{Oh})$ at high $l$
sets in almost immediately away from any point in the operative range, so no
placement of a single anchor buys a genuinely wide small-parameter regime.
\end{derivbox}

\begin{keybox}[The honest conclusion]
For this fluid's fitted parameters, there is no single reference viscosity
-- rest, infinite-shear, or any point in between -- around which a
regular (or matched-asymptotic, single-anchor) perturbation expansion of
Reid's damping rate is uniformly valid across the modes this solver
resolves. The physics genuinely lives in the fully nonlinear, transitional
part of the Carreau-Yasuda curve, not in a neighborhood of either plateau.
This is a quantitative fact about \emph{this fluid's} $\lambda_c$ and $a$,
not a universal statement about shear-thinning drops.
\end{keybox}

%% ────────────────────────────────────────────────────────────────────────────
\subsection{What Reid's problem already gives you for free}
\label{sec:cy_exact_evaluation}
%% ────────────────────────────────────────────────────────────────────────────

Because the failure above is specifically a failure of \emph{linearizing}
$\lambda_l(\mathrm{Oh})$, $\omega_l^2(\mathrm{Oh})$, and because these
functions are already known \emph{exactly}, for any $\mathrm{Oh}$, via
Reid's characteristic equation (solved numerically by continuation, and
tabulated for fast repeated evaluation), the correct response is not to
build a better linear correction -- it is to not linearize this piece at
all. Given an instantaneous effective Ohnesorge number
$\mathrm{Oh}_{\mathrm{eff}}(t)$ from any closure (including the one derived
in Section~\ref{sec:cy_temporal} below), the exact, non-perturbative
evaluation of $\lambda_l(\mathrm{Oh}_{\mathrm{eff}}(t))$,
$\omega_l^2(\mathrm{Oh}_{\mathrm{eff}}(t))$ has no validity-range caveat and
is already the practice this repo's existing (heuristic) Carreau-Yasuda
extension follows. That part of the existing scheme was never the weak
point. The weak point is upstream of it: whether it is valid to replace a
velocity field's genuinely \emph{spatially varying} local viscosity with a
single scalar $\mathrm{Oh}_{\mathrm{eff}}(t)$ fed into a formula built for a
spatially \emph{uniform} viscosity. That is the question the rest of this
section answers.

%% ────────────────────────────────────────────────────────────────────────────
\subsection{Machinery I: adjoint sensitivity for Reid's boundary-value problem}
\label{sec:cy_adjoint}
%% ────────────────────────────────────────────────────────────────────────────

Reid's velocity equation, $\mathcal{L}[U]\equiv U''-l(l+1)U/x^2+q^2U$
(Section~\ref{sec:velocity_ode}), is formally self-adjoint on $(0,1)$ with
the plain $L^2$ inner product -- no weight function is needed, because
$\mathcal{L}$ has no first-derivative term:
\begin{derivbox}[Lagrange identity for $\mathcal{L}$]
For any twice-differentiable $U,W$,
\begin{equation}
  W\,\mathcal{L}[U] - U\,\mathcal{L}[W]
  = \frac{d}{dx}\Bigl[\,W\,U' - W'\,U\,\Bigr].
  \label{eq:lagrange}
\end{equation}
Verified symbolically (Symbolics.jl, abstract $U(x)$, $W(x)$): both sides
agree identically.
\end{derivbox}

This gives a shortcut for exactly the calculation a perturbation of Reid's
problem needs: the boundary derivative of a particular solution, without
solving the ODE.

\begin{keybox}[Adjoint shortcut]
Let $Y(x)$ be regular at $x=0$, satisfy $Y(1)=0$, and solve
$\mathcal{L}[Y]=\mathrm{RHS}(x)$ for a known forcing $\mathrm{RHS}$, at a
fixed wavenumber $q_0$ (so $\mathcal{L}$ uses $q_0^2$). Then
\begin{equation}
  Y'(1) = \frac{1}{j_l(q_0)}\int_0^1 x\,j_l(q_0x)\,\mathrm{RHS}(x)\,dx,
  \label{eq:adjoint_shortcut}
\end{equation}
using the regular homogeneous solution $x\,j_l(q_0x)$ as the adjoint test
function (both boundary terms at $x=0$ vanish by regularity; the boundary
term at $x=1$ collapses using $Y(1)=0$).
\end{keybox}

\begin{notebox}[Numerically verified, not just derived]
Checked against direct numerical shooting (regular-solution seeding near
$x=0$, RK4 integration, homogeneous-solution correction to enforce
$Y(1)=0$) for three independent $(l,q_0,\mathrm{RHS})$ combinations: relative
agreement of $10^{-9}$ to $10^{-11}$ once the shooting method's own
near-origin seed is chosen at a moderate radius ($x_0\sim 10^{-2}$, not
$10^{-6}$ -- the regular solution scales as $x^{l+1}$ near the origin, and
seeding too close to $0$ loses precision to floating-point underflow before
the adjoint comparison is even meaningful).
\end{notebox}

Reid's boundary conditions close the system: BC1 ($U(1)=-1$) is unaffected
by any viscosity correction, since it is purely kinematic, and BC2
($\mathcal L_2[U](1)=0$ plus a known correction) is linear in the same
unknowns the shortcut already produces. Combining the two turns the search
for the $O(\delta)$ shift in $q^2$, from \emph{any} new forcing added to
Reid's ODE and both stress boundary conditions, into an ordinary linear
solve -- no further ODE integration is required. This reduction has been
worked through by hand for the general case (not yet promoted to an
assertion in the backing script, since the resulting expression is a large
symbolic fraction rather than a compact closed form); what \emph{is}
CAS-verified end to end is the shortcut itself
(Eq.~\ref{eq:adjoint_shortcut}), which is the genuinely nontrivial step.

%% ────────────────────────────────────────────────────────────────────────────
\subsection{Machinery II: the strain-rate field of Reid's actual viscous mode}
\label{sec:cy_strain}
%% ────────────────────────────────────────────────────────────────────────────

Any first-principles shear-rate calculation must be built from Reid's actual
\emph{viscous} velocity profile $U(x)=C\,x\,j_l(qx)+\Pi_0x^{l+1}$
(Section~\ref{sec:velocity_ode}), not the inviscid potential-flow shape
$r^lP_l(\cos\theta)$ that this repo's existing heuristic scripts use to
estimate a characteristic shear rate. The two are not interchangeable: Reid's
own damping normalization comes from the homogeneous (Bessel) part of
$U(x)$, which \emph{is} the viscous correction to potential flow -- using
the potential-flow shape to estimate shear rate is inconsistent with Reid's
theory already at zeroth order, before any shear-thinning correction is
applied.

Writing $u_r=F(x)P_l(\cos\theta)$, $F(x)\equiv U(x)/x^2$, and $u_\theta$
from the stream-function relation of Section~\ref{sec:bc2}
($u_\theta=-(2F+xF')\sin\theta\,P_l'(\cos\theta)/[l(l+1)]$), the standard
axisymmetric strain-rate tensor is
\begin{align}
  e_{rr}   &= F'(x)\,P_l(\mu), \\
  e_{\theta\theta} &= \frac{1}{x}\frac{\partial u_\theta}{\partial\theta} + \frac{u_r}{x}, \\
  e_{\varphi\varphi} &= \frac{u_r + u_\theta\cot\theta}{x}, \\
  e_{r\theta} &= \frac12\left[x\frac{\partial}{\partial x}\!\left(\frac{u_\theta}{x}\right) + \frac{1}{x}\frac{\partial u_r}{\partial\theta}\right],
  \qquad \mu\equiv\cos\theta.
  \label{eq:strain_tensor}
\end{align}

\begin{notebox}[Incompressibility check]
$e_{rr}+e_{\theta\theta}+e_{\varphi\varphi}=0$ identically, for \emph{any}
radial profile $F(x)$ -- verified both symbolically (concrete Legendre
polynomials, $l=2,3$, after resolving a $\sin^2\theta+\cos^2\theta=1$
simplification Symbolics.jl does not apply automatically) and numerically
(concrete polynomial $F$, $l=2,\dots,8$, to floating-point precision at
every tested $(x,\theta)$). This is a kinematic identity of the poloidal
representation itself, confirming the strain-tensor construction is
consistent regardless of what $F(x)$ turns out to be.
\end{notebox}

%% ────────────────────────────────────────────────────────────────────────────
\subsection{Machinery III: the period-$\pi$ lemma}
\label{sec:cy_period_pi}
%% ────────────────────────────────────────────────────────────────────────────

At $\mathrm{Oh}_\infty$, $q$ is genuinely complex ($q_0\approx
7.60-7.30\,\mathrm{i}$ for $l=2$), so Reid's mode is a true damped
oscillation, not a pure decay. The physical (real) strain field at phase
$\phi\equiv\omega t$ is $\mathrm{Re}\bigl[e_{ij}(x,\theta)\,e^{-\mathrm{i}\phi}\bigr]$
for each complex tensor component $e_{ij}$ built from~\eqref{eq:strain_tensor}
with complex $F$. A natural first attempt at a shear-thinning correction --
evaluate everything at one representative phase, e.g. $\phi=0$ -- turns out
to be indefensible: numerically, the real physical shape at $x_0=0.7$
swings from $-0.345$ to $+0.345$ across one period, nowhere near constant.
Any nonlinear (in particular, fractional-power) function of this field
therefore depends on \emph{which} phase was chosen, unless the choice is
justified by an actual time-average.

\begin{keybox}[The shear-rate invariant cannot resonate at the mode's own frequency]
Let $S\equiv\sqrt{2\,e_{ij}e_{ij}}$ (summed over the strain-tensor
components) be built from any single-frequency oscillating axisymmetric
poloidal field, real or complex amplitude. Then $S(\phi)$ is exactly
periodic with period $\pi$ in $\phi=\omega t$, not $2\pi$, and hence its
Fourier series over the full period contains only even harmonics: the
coefficient of $\cos(\phi)$ and $\sin(\phi)$ (the fundamental, $m=1$) is
identically zero, for $S$ itself and for \emph{any} function of $S$
(fractional power or otherwise).
\end{keybox}

\begin{derivbox}[Proof]
Each strain component satisfies
$\mathrm{Re}(e_{ij}e^{-\mathrm{i}\phi})=|e_{ij}|\cos(\phi-\arg e_{ij})$.
Squaring, $|e_{ij}|^2\cos^2(\phi-\arg e_{ij}) =
\tfrac12|e_{ij}|^2\bigl[1+\cos(2\phi-2\arg e_{ij})\bigr]$, which is
manifestly invariant under $\phi\to\phi+\pi$. $S^2$ is a sum of such terms,
hence itself period-$\pi$; $S=\sqrt{S^2}\ge0$ and any real function $h(S)$
inherit the period exactly (no branch ambiguity, since $S^2\ge 0$
throughout). A period-$\pi$ function integrated against $\sin(m\phi)$ or
$\cos(m\phi)$ over $[0,2\pi)$ vanishes identically for every odd $m$, since
the two half-period copies enter with opposite sign and cancel.
\end{derivbox}

\begin{notebox}[Confirmed two ways]
Numerically: $S^2(\phi+\pi)=S^2(\phi)$ to $10^{-15}$--$10^{-16}$ at multiple
independent $(x,\theta)$ points. Via direct Fourier projection on a joint
$800\times800$ $(\theta,\phi)$ grid: the $m=1$ coefficient (both $\cos$ and
$\sin$) is zero to $10^{-17}$--$10^{-18}$ at every spherical-harmonic degree
$l'$ tested ($0,2,4,6,8$), while $m=0$ and $m=2$ are the genuinely nonzero
temporal channels.
\end{notebox}

This has an immediate structural consequence: since no forcing term
resonant at the base mode's own frequency exists, the adjoint machinery of
Section~\ref{sec:cy_adjoint} -- built for exactly this kind of resonant
solvability condition -- has nothing to act on at $m=1$. The leading
temporal effect of any generalized-Newtonian correction on this mode is
therefore carried entirely by the $m=0$ (period-averaged, effectively
steady) channel; the $m=2$ channel represents a distinct, second-harmonic
\emph{parametric} coupling (a $2\!:\!1$ resonance structure with the base
oscillation), a different physical phenomenon not addressed here.

%% ────────────────────────────────────────────────────────────────────────────
\subsection{The genuinely open question: spatial homogenization}
\label{sec:cy_temporal}
%% ────────────────────────────────────────────────────────────────────────────

Because the $m=0$ channel is the leading one, and because the adjoint
shortcut of Section~\ref{sec:cy_adjoint} places no restriction on the
forcing's spatial shape -- a steady (non-oscillating) correction is exactly
as valid an input to it as any other -- the machinery built here is a
reasonable tool to bring to the question that actually remains open (this
specific reduction has not itself been carried through to a checked
assertion, and is flagged as such rather than presented as settled): replacing a
spatially varying $\eta(x,\theta)$ with a single scalar
$\mathrm{Oh}_{\mathrm{eff}}(t)$ is an uncontrolled effective-medium
approximation unless the spatial variation is itself small. Checked
directly, using the period-averaged shear-rate shape from
Section~\ref{sec:cy_strain}--\ref{sec:cy_period_pi} over the drop's volume
at $\mathrm{Oh}_\infty$, $l=2$: the shape varies by a factor of
$\approx 2.25$ between its minimum and maximum over $(x,\theta)$ ($-50\%$ to
$+13\%$ around the mean). Propagated through the Carreau-Yasuda exponent
$a\approx0.743$, this is roughly a factor-of-two spread in local viscosity
across the drop at a single instant -- smaller than the multiple-order-of-
magnitude range $\mathrm{Oh}_{\mathrm{eff}}$ traverses over an impact
\emph{in time}, but not small in the sense a first-order correction needs
to be uncontroversial.

Legendre-projecting the (properly $m=0$, period-averaged, not single-phase)
shear-rate correction shape at a representative radius confirms the
fractional-power leakage mechanism anticipated on general grounds: a
\emph{single} active mode ($l=2$) spontaneously forces every even
spherical-harmonic degree ($l'=0,4,6,8$ all genuinely nonzero, odd $l'$
zero to $10^{-11}$ by a verified $\cos\theta$-parity argument), because
raising a finite Legendre polynomial to a non-integer power does not, in
general, stay within its own degree.

\begin{keybox}[What this section establishes, precisely]
\begin{enumerate}
  \item The classical $O(\epsilon^3)$ Carreau expansion of
        Section~\ref{sec:carreau} does not extend to this fluid's fitted,
        non-integer Carreau-Yasuda exponent: no small-amplitude anchor,
        at rest, at infinite shear, or in between, supports a uniformly
        valid linear correction to $\lambda_l(\mathrm{Oh})$ across the
        modes this solver resolves.
  \item Reid's exact (non-perturbative, continuation-solved)
        $\lambda_l(\mathrm{Oh})$, $\omega_l^2(\mathrm{Oh})$ should be
        evaluated directly at whatever $\mathrm{Oh}_{\mathrm{eff}}(t)$ a
        closure provides -- never linearized -- since this repo already
        implements a fast, tabulated version of exactly that evaluation.
  \item The genuinely open approximation in the existing heuristic scheme
        is the collapse of a spatially varying $\eta(x,\theta,t)$ onto a
        single scalar $\mathrm{Oh}_{\mathrm{eff}}(t)$, quantified here at
        roughly a factor of two in local viscosity across the drop at a
        representative instant -- a real, moderate, but not yet
        controlled error.
  \item Three general, CAS-verified results survive independently of any
        of the above: the adjoint sensitivity shortcut
        (Eq.~\ref{eq:adjoint_shortcut}), the viscous strain-rate tensor
        (Eq.~\ref{eq:strain_tensor}), and the period-$\pi$ lemma. None
        depend on this fluid's specific parameters, and all three are the
        correct starting point for a future, genuinely first-principles
        treatment -- most plausibly a self-consistent, spatially resolved
        solve (using the adjoint machinery as a Newton/Fr\'echet-derivative
        iteration step, not a one-shot small-parameter expansion) rather
        than any further attempt at a closed-form amplitude equation.
\end{enumerate}
\end{keybox}

\begin{notebox}[Scope]
This section does not implement the self-consistent spatially resolved
solve it recommends -- that is future work. It also does not revisit the
$m=2$ parametric channel identified in
Section~\ref{sec:cy_period_pi}, nor extend the strain-tensor/leakage
calculation beyond $l=2$ and a single representative radius. What it does
establish, to the same CAS-verified standard as every other section of this
document -- every numbered claim above has a corresponding assertion in
\texttt{julia/derivations/carreauYasuda\_firstprinciples\_derivation.jl},
including a live cross-check against a running \texttt{solve\_drop!} trace
(the pinned-intermediate-Oh finding of
Section~\ref{sec:cy_anchors}) -- is precisely which of the natural next
steps work, which do not, and why -- for this fluid, not as a general
verdict on Carreau-Yasuda drops.
\end{notebox}
